\newcommand{\wineIntro}{But tonight, we also reflect not just on our desires to bring freedom, but on concrete strategies that may help us be more effective in doing good in the world.}

\newcommand{\wineSummaryThisYear}{This year, we learn about tools to be more effective at bringing freedom to ourselves and our world.}

\newcommand{\wineFirstCup}{
    \newReading With this first cup of wine, we consider the philosophy of ``Effective Altruism''”'' as a way to bring greater freedom to the oppressed.

    Reducing the number of children exposed to gang violence in Latin America, bridging the education gap in America, or eradicating AIDS are all worthy goals that would bring positive change to the world. But these lofty goals are too difficult to solve in an evening, no matter how good our intentions. 

    Just because something is difficult does not mean it is not worthwhile. But for this Seder we focus on something more directly at hand\textemdash{}something we each can accomplish as individuals in the next few weeks.
    This is not to forget the overarching goals of \cjRL{tiqUn `OlAm} (\textit{tea-koon oh-lawm})\textemdash{}``repairing the world'' and bringing freedom to the oppressed\textemdash{}but to look, as a first step, to the views Australian philosopher Peter Singer describes as ``Effective Altruism''. The Effective Altruism philosophy can be summarized as ``work smarter, not harder to maximize your positive impact in social justice.'' During this Seder we investigate one philosophy of how to train ourselves to be more effective agents for bringing freedom to our world. 

    \subsubsection*{First World Freedom}
    \newReading In this country today, the common man takes for granted more abundance and pleasure than what was afforded to even the kings of yesteryear. King Solomon's table did not contain even half the selection of worldwide foods found at a mall's food court or your local supermarket.

    We are free to learn; we carry all of humanity's knowledge with us in our pocket, accessible at our fingertips. Anytime we please, we can sit at the virtual feet of the world's wisest sages listening to their TED Talks.

    We are free to travel not only our country, but the entire planet. A mere two hundred years ago, people sold themselves into seven years of indentured servitude to pay for passage across the Atlantic Ocean, whereas today one can purchase airfare with fewer than seven days salary, even at the federal minimum wage. 

    \newReading And yet, despite these apparent freedoms and luxuries, nearly everyone feels that they are not truly free. That is, they feel they \textit{do not have the opportunity to live life as they want to live it}.

    Striving for such a lofty definition of freedom may seem like the ultimate ``first world problem.'' But if freeing the students trapped in a struggling school system is your goal, you will likely be better equipped and more effective in achieving that goal after you have completely freed yourself. As the astute flight attendants point out: ``Secure your own mask first before helping others.''

    This Seder we will take a journey adapted from Harry Browne's book \textit{How I Found Freedom in an Unfree World} to discover why we do not feel free, what keeps us from feeling free, and how we can become much more free. 

    With this first cup of wine, we think about the idea that by first freeing ourselves, we become more effective in our quest to bring freedom to others.
}

\newcommand{\wineSecondCup}{
    \newReading With this second cup of wine, we contemplate why we\textemdash{}even with the wealth and priveleges we hold\textemdash{}do not feel truly free.
    
    \subsubsection*{Freedom}
    To understand why we do not feel free as individuals\textemdash{}even in a society that by objective measures enjoys the most freedom the world has ever seen\textemdash{}we must first better understand how we as members of modern Western society perceive freedom.

    For most people, freedom is an ``if only.'' ``If only it hadn't been for my husband / mother / children / boss, I would have been a success.'' Or ``If only it hadn't been for Carter (or Bush, or Obama, or whomever), the country would be free.''

    No matter which specific freedom you crave most\textemdash{}freedom from social restrictions, family problems, high taxes, bad relationships, the rat race of work, governmental repression\textemdash{}the first step on the journey to feeling free is to choose to live that way. Concentrate on the things you control, and use that control to remove the restrictions and complications from your life.

    Tonight, together, we will explore ways to move down the path towards greater freedom.

    \subsubsection*{Happiness}
    \newReading Ultimate freedom lies in the ability to pursue happiness, and we feel unfree when we think someone is preventing us from achieving happiness.

    In everything you do, with the knowledge and insight at your disposal, you choose what you think will give you the most well-being and the least mental discomfort. The objective is what is usually called \textit{happiness}\textemdash{}the feeling of well-being.
    Happiness isn't a new car, fame, a good marriage, wealth, or a warm blanket. Those are things. Happiness is what you feel inside of you as a result of the things that happen to you.

    Happiness might be \textit{produced} by a good marriage, fame, a new car, or a warm blanket. For some people, happiness occurs as a result of doing favors for other people; for others, it results from bringing about social reforms; for still others, it comes from believing they've outsmarted someone. It might come from a big meal, music, dancing, singing, working, kissing, resting, feeding the homeless, mentoring at-risk youth, etc. These are things that might make an individual feel good.

    \subsubsection*{Decisions}
    \newReading Everything you do is motivated by the desire to feel as much happiness as possible and to eliminate mental discomfort\textemdash{}either in the short term or the long term.
    For example, you may work hard at your career for many years because you feel ``it will all be worth it someday''\textemdash{}meaning it will enable you to do the things that will make you feel good. Or perhaps the job itself makes you feel better than anything else you've considered.
    Or you may do certain things because you're afraid that if you don't do them you'll feel bad. You may lend money to your friends only because you'd feel guilty if you didn't.

    A \textit{positive decision} is one in which you choose among alternatives to maximize your happiness. An example would be deciding whether you'll be happier going to a movie or a football game.
    A \textit{negative decision} is one in which you choose among alternatives to minimize your unhappiness. An example would be deciding whether to let your roof leak or to empty your savings account to get it fixed. Neither choice will increase your happiness; you're trying to decide which choice would be the least unpleasant.

    \newReading A free person spends most of his time making positive decisions\textemdash{}choosing among attractive alternatives.

    Deciding to spend your time helping the less fortunate is almost always made as a positive decision; you have decided that doing so will make you happier. People making negative decisions are usually\textemdash{}by definition\textemdash{}trying to solve their own problems. Therefore, to maximize your ability to help others, you must first free yourself to be able to spend your time making positive decisions instead of negative ones.

    We drink this second cup with the hope that things we discuss later tonight might help us move to a freer life where we spend more of our time making positive choices rather than negative ones. 
}

\newcommand{\wineThirdCup}{
    \newReading With the third cup of wine, we think about a common belief that restricts our freedom and how to overcome it.

    \subsubsection*{Traps}
    Our culture is saturated with philosophical ``truths'' that are commonly accepted and acted upon\textemdash{}and are rarely challenged. Let us think of these truisms as \textit{traps}. It's very easy to get caught in a trap. The truisms are repeated so often they can be taken for granted. And that can lead to acting upon the suggestions implied in them\textemdash{}resulting in wasted time, fighting inappropriate battles, and attempting to do the impossible.

    If you do not feel free now, it's very likely that you've accepted some of these traps. And you probably haven't known of a number of alternatives that could get you out of your restrictions without the pain and effort you might have assumed would be necessary.
    As we look at these traps and alternatives, try and become aware of the unlimited number of avenues open to you. You possess a tremendous amount of control over your situation\textemdash{}control that's disregarded when you focus attention upon the people who seem to stand in your way.

    \subsubsection*{The Identity Trap}
    \newReading There are two forms of the Identity Trap: (1) the belief that you should be someone other than yourself; and (2) the assumption that others will do things in the way you would.

    In the first trap, you necessarily forfeit your freedom by requiring yourself to live in a stereotyped, predetermined way that doesn't consider your own desires, feelings, and objectives. The second trap is more subtle, but just as harmful to your freedom. When you expect someone to have the same ideas, attitudes, and feelings you have, you expect them to act in ways that aren't in keeping with their nature. As a result, you'll expect and hope that people will do things they're not capable of doing.

    Let's begin by recognizing what we know about you. We know \textit{you are different}. You're different from everyone else in the world.

    \newReading Just as no two persons' fingerprints are identical, no two people are identical in terms of their knowledge, understanding, attitudes, likes, and dislikes. No one else has lived your life and experienced all the same things. What you consider to be logic or common sense will vary in some way from another person's logic. As a result, you see and interpret and react to what goes on around you differently from anyone else.

    That's not hard to see. And yet the root of the Identity Trap is that most individuals are oblivious to these differences. They assume that all people want the same things\textemdash{}or that they \textit{should} want them. They expect everyone to respond in the same way to the same things. They assume that what one person or elite group sees and accepts should be accepted by everyone. This can range from one person expecting another to enjoy the movie he enjoys to the individual who's upset if everyone doesn't go to synagogue.

    You are you and only you. You live in a world of your own, composed of your own experience. You can't be someone other than who you are, and you are the most qualified person in the entire world to determine what will make you feel happy and free.

    \subsubsection*{Avoid Trapping Yourself}
    \newReading You’re in the Identity Trap when you try to deny your bad feelings\textemdash{}such as hate, fear, jealousy, or guilt. You're in the trap when you try to deny good feelings\textemdash{}such as enjoyment of something that's frowned upon. You're in the trap when you try to make yourself feel good about someone or something that you don't enjoy\textemdash{}such as going to see an opera just to prove you are ``cultured''.

    You're also in the trap when you believe you should be happy simply because you're doing what you've been told will make you happy. Take the example of the businesswoman who has to keep reminding herself that her six-figure salary and carpeted office are what she's always wanted. Or the man who keeps telling himself he must be happy, now that he finally has a wife, four children, and a home in suburbia. Each of them is living a life they've been told should make them happy; but if it doesn't, they attempt to make their emotions respond.

    \newReading To escape the trap you must avoid automatically conforming to what you have heard will make other people happy, and instead understand if what you are doing is on the path that will make \textit{you} happy.

    One strategy is to go through a whole week (or maybe just a whole day to start) and take five minutes each hour to think about the past fifty-five; write down at least one thing you did that truly made you happy, and try to uncover an action you took because you were trapped into thinking that it \textit{should} make you happy even though it doesn't.

    The use of this technique can dramatically increase your self-control, self-confidence, and happiness. It will help provide an answer to the important question:
How many hours a week are you happy?

    \subsubsection*{Avoid Letting Others Trap You}
    \newReading You're also in the Identity Trap when you assume an individual will react to something as you would react or as you've seen someone else react. You could make everyone else be, act, and think in ways of your choosing if you were \god{}. But you aren't. So it's far more useful to recognize and accept each person as he is\textemdash{}and then deal with them accordingly. \textit{You can't control the natures of other people, but you can control how you'll deal with them}. And you can also control the extent and manner in which you'll be involved with them.

    The paradox is that you have tremendous control over your life, but you give up that control when you try to control others.

    To escape this form of the Identity trap, you must understand the concept of \textit{direct} and \textit{indirect} alternatives. A \textit{direct} alternative is one that requires only direct action \textit{by yourself} to get a desired result. An \textit{indirect} alternative requires that you act to make \textit{someone else} do what is necessary to achieve your objective. 

    \newReading For example, suppose you're a college student who's dissatisfied with the curriculum offered at your school. Examples of indirect alternatives would be to try to influence the administration to change the curriculum or try to influence higher powers to change the administration. Examples of direct alternatives would be to find a more agreeable school or to study the missing subjects on the side.

    A direct alternative requires only a decision on your part; an indirect alternative requires that you get others to make decisions that would benefit you. 

    With this third cup of wine, we think about how we can escape the Identity Trap by using direct alternatives to more easily solve most situations, thereby freeing us to have the time and energy to attempt indirect alternatives only for more intractable causes and goals that would genuinely bring us happiness to achieve.
}

\newcommand{\wineFourthCup}{
    \newReading As our Seder draws near to the end, we take up our cups one last time as we think about how we can further accelerate the exponential growth to freedom.
    
    \subsubsection*{Driving Innovation}
    While universal freedom is a very real possibility, we're also in a race against time. Until the innovations of abundance bear fruit, scarcity remains a real concern. And nearly as bad as scarcity is the threat of scarcity and the devastating violence it can often incite. How can we steer innovation and step on the gas?
    
    Decades ago, Buckminster Fuller observed that ``You never change things by fighting the existing reality. To change something, build a new model that makes the existing model obsolete.'' In our current era of exponential technological growth, where we are so close on so many fronts to achieving universal freedom, many of the most high-profile philanthropists have taken Fuller's quotation to heart and changed how they donate their money. Rather than solely\textemdash{}and sometimes even instead of\textemdash{}supporting traditional charities or NGOs, they are increasingly supporting organizations that spur radical innovation aimed at addressing societal problems.
    
    \newReading The most famous of these organizations is the X PRIZE Foundation, which launches challenges offering large cash prizes to any team who can solve a pressing problem; automobile fuel efficiency, global learning, women's safety, and atmospheric water extraction are just some of the goals they have set. One of the most successful of these challenges was launched in the wake of the \textit{Deepwater Horizon} BP oil spill. Shocked at the magnitude of environmental damage that occurred, philanthropist Wendy Schmidt agreed to bankroll the competition. The focus of the prize was clear: the technology used to clean up the BP spill in 2010 was the same as was used to clean up the Exxon Valdez spill in 1989\textemdash{}it was clearly time for an upgrade. Less than 18 months after the BP oil spill, a team won the challenge with a new process that had quadrupled the performance of the oil industry's existing cleanup technology. What made this even more exciting was that this giant breakthrough had been achieved just by offering a \$1 million prize. Considering that the environmental damage of just that single spill was estimated at \$17 billion, the ability to prevent environmental damage by cleaning up all future spills more effectively made the return on charitable investment enormous.
    
    People across the spectrum are getting behind the concept of incentive prizes as an extremely effective method of achieving social change. Software entrepreneurs such as Elon Musk, Larry Page, and Sergey Brin, motivational speaker Tony Robbins, film director James Cameron, columnist Arianna Huffington, and former first lady Barbara Bush are just a few of the strong supporters of the X PRIZE.
    
    \newReading Frustrated after an oil company managed to secure drilling rights to land in the Amazon Rainforest\textemdash{}that a charity he had donated to was supposed to protect\textemdash{}philanthropist Vishen Lakhiani switched his support to the X PRIZE Foundation and explained his reasoning in a company memo: ``Why keep fighting companies acre by acre through the rainforest when we can work to make oil unnecessary with solar power innovation? Or make food production so efficient that rainforest doesn't need to be cut down.''

    You don't have to be rich enough to sponsor your own X PRIZE to make a difference. Anyone can decide to donate to support part of the next great humanitarian breakthrough.
    
    \subsubsection*{Don't Fall Victim to Cognitive Biases}
    The more people who understand the unprecedented opportunity for freedom that the near future holds, the faster we'll achieve this age-old goal. It can be tempting to dismiss all these thoughts and ideas about the potential for universal freedom and abundance as na\"ive optimism; our brains are wired to jump to that conclusion. But don't let your biology control you. Recognize these biases that are luring you to reject this exciting future, and point them out to others who fall into their trap.
    
    Confirmation bias is a tendency to search for or interpret information in a way that confirms one's preconceptions\textemdash{}but it can often limit our ability to take in new data and change old opinions. This bias is typified in the huge uptick in downright lies told by politicians that still persist in large portions of the public's minds because people are able to find so many other ``sources'' on the internet or cable news networks that repeat that false information. If you are thinking that there's no way we can free the poor because the hole we're in is too deep to climb out of, check that you are actually giving new information the same weight as what you already ``know.''
    
    \newReading Anchoring bias is the predilection for relying too heavily on one piece of information when making decisions. When people believe the world's falling apart, it's often an anchoring problem. At the end of the 19\textsuperscript{th} century, London was becoming uninhabitable because of the accumulation of horse manure. People were absolutely panicked. Because of anchoring, they couldn't imagine any other possible solutions. No one had any idea the car was coming and soon they'd be worrying about dirty skies, not dirty streets.
    
    Just because current water purification technology requires too much electricity to deploy effectively in developing countries shouldn't make you forget that abundant energy is coming in the next few years that will completely alter those costs.
    
    The negativity bias\textemdash{}the tendency to give more weight to negative information and experiences than positive ones\textemdash{}sure isn't helping matters either. That bias gets further compounded by the strategy of news sources to follow the ``if it bleeds, it leads'' philosophy of almost exclusively reporting violent and negative news because it gets higher ratings; often over 90\% of articles are pessimistic.
    
    \newReading Knowing about these biases helps keep you on guard and better able to seek out the truth. Finally, if you ever find yourself really believing that the world is actually getting worse, remember these thoughts from evolutionary psychologist Steven Pinker:\par
    \noindent\makebox[\textwidth][c]{
        \begin{minipage}[t]{0.9\textwidth}
            Cruelty as entertainment, human sacrifice to indulge superstition, slavery as a labor-saving device, conquest as the mission statement of government, genocide as a means of acquiring real estate, torture and mutilation as routine punishment, the death penalty for misdemeanors and differences of opinion, assassination as the mechanism of political succession, rape as the spoils of war, pogroms as outlets for frustration, homicide as the major form of conflict resolution\textemdash{}all were unexceptional features of life for most of human history. But, today, they are rare to nonexistent in the West, far less common elsewhere than they used to be, concealed when they do occur, and widely condemned when they are brought to light.
        \end{minipage}
    }
    
    With this fourth cup of wine, we understand that bringing freedom to all the world is not some abstract, never ending task\textemdash{}as it was during every past Seder since we left Egypt\textemdash{}but something actually achievable in our lifetimes. Use that as motivation, spread the good news to other people, and work to make universal freedom happen even sooner than predicted.


}